\chapter{Qualit\`a del software e verifica}
\label{cap:26-qualita-verifica}

In un progetto che si propone come strumento di misura, la verifica non \`e
igiene ma parte del contributo: un risultato ottenuto su una piattaforma di cui
non si conosce lo stato non \`e un risultato. Questo capitolo descrive
l'apparato di verifica di \nf{}, dalle unit\`a alle regole architetturali.

\section{Le proporzioni}

\begin{center}
\small
\begin{adjustbox}{max width=\textwidth}
\begin{tabular}{@{}lr@{}}
\toprule
\textbf{Grandezza} & \textbf{Valore} \\
\midrule
Righe di codice di produzione (Java) & 16\,487 \\
Righe di codice di test (Java) & 31\,664 \\
Rapporto test / produzione & $\approx 1{,}9$ \\
Classi di test & 235 \\
Soglia minima di copertura (JaCoCo) & 85\,\% \\
\bottomrule
\end{tabular}
\end{adjustbox}
\end{center}

Quasi due righe di test per ogni riga di produzione. Il rapporto \`e alto ma non
sorprendente per un progetto la cui superficie \`e piccola e le cui interazioni
--- fra moduli, fra code, fra politiche --- sono numerose: il codice combinatorio
\`e nei test, non nella produzione.

\section{I livelli}

\begin{center}
\small
\begin{adjustbox}{max width=\textwidth}
\begin{tabular}{@{}lp{8.4cm}@{}}
\toprule
\textbf{Livello} & \textbf{Caratteristiche} \\
\midrule
Unit\`a & JUnit 5, classi \code{*Test.java}, nessuna infrastruttura. \\
Integrazione & Contesto Spring completo con \code{WebTestClient}; server mock
Fabric8 per Kubernetes. \\
Container & Testcontainers~\cite{testcontainers}: richiede Docker. \\
Kubernetes & Richiede \code{KUBECONFIG}. \\
Contratto multi-SDK & \code{functions/contract-tests/run.sh}
(\cref{cap:18-sdk}). \\
Watchdog & Test unitari Rust pi\`u integrazione locale per HTTP, STDIO, FILE,
callback, metriche e ciclo di vita warm. \\
E2E di scenario & Di propriet\`a dello strumento esterno
(\cref{cap:22-metodo-sperimentale}). \\
\bottomrule
\end{tabular}
\end{adjustbox}
\end{center}

Va notato che i test del watchdog non richiedono n\'e Docker n\'e una VM: bastano
Rust, Python 3 e alcune utilit\`a a riga di comando. \`E una scelta che mantiene
verificabile in locale il componente presente in ogni container di funzione.

\section{Le due invocazioni della suite}

Come discusso nel \cref{cap:16-moduli-diagnostici}, la selezione modulare \`e
risolta a tempo di configurazione di Gradle e non pu\`o variare all'interno di
una build. La verifica completa richiede quindi due invocazioni:

\begin{lstlisting}[language=bash]
./gradlew test                              # default: tutti i moduli
./gradlew test -PcontrolPlaneModules=none   # la configurazione minima
\end{lstlisting}

\section{Le regole architetturali}

Il progetto codifica sei regole con ArchUnit~\cite{archunit}. Cinque riguardano
il nucleo e una i moduli.

\subsection{Le regole del nucleo}

\begin{description}[style=nextline, leftmargin=0pt]
  \item[\textbf{R1 --- Assenza di cicli.}] I package di primo livello del nucleo
  non devono formare cicli fra loro.
\begin{lstlisting}[language=Java]
slices().matching("..controlplane.(*)..").should().beFreeOfCycles();
\end{lstlisting}

  \item[\textbf{R2 --- \code{api} \`e l'unico punto di ingresso.}] Nessuna classe
  fuori dal package \code{api} deve dipendere da esso. La direzione \`e
  unidirezionale: l'API conosce il servizio, il servizio non conosce l'API.

  \item[\textbf{R3 --- Il nucleo non dipende dai moduli.}]
\begin{lstlisting}[language=Java]
noClasses().that().resideInAPackage("..controlplane..")
        .should().dependOnClassesThat().resideInAPackage("..modules..")
        .as("the core must work without optional modules (SPI)");
\end{lstlisting}
  \`E la regola che rende \emph{verificata} la promessa del
  \cref{cap:03-requisiti-principi}. Senza di essa, l'affermazione ``il nucleo
  funziona senza moduli'' sarebbe un'intenzione; con essa \`e una propriet\`a che
  la build accerta.

  \item[\textbf{R4 --- La direzione dei livelli.}] I package \code{dispatch},
  \code{execution} e \code{deployment} non devono dipendere da \code{service}. La
  gerarchia \`e \code{api} $\to$ \code{service} $\to$ \code{dispatch}, e i livelli
  bassi non risalgono.

  \item[\textbf{R6 --- Il namespace appartiene al nucleo.}] Ogni classe nel
  package \code{it.unimib.datai.nanofaas.controlplane} deve avere il proprio
  file di classe sotto \code{platform/control-plane/}.
\end{description}

L'ultima \`e la pi\`u insolita e vale la pena soffermarcisi. Verifica non una
dipendenza ma una \emph{corrispondenza fra nome logico e collocazione fisica}:
impedisce che un modulo dichiari una classe dentro il namespace del nucleo per
aggirare le altre regole. Senza di essa, R3 sarebbe eludibile spostando il
package anzich\'e la dipendenza.

L'implementazione ha dovuto adattarsi a un cambiamento della libreria:

\begin{lstlisting}[language=Java]
// ArchUnit 1.4.1 dropped SourceCodeLocation's source-path accessor, so read the
// source URI instead; its path points under platform/control-plane/.
String path = item.getSource().map(Source::getUri).map(URI::getPath).orElse("");
\end{lstlisting}

\subsection{La regola dei moduli}

Ciascuno dei nove moduli porta la propria classe di verifica, con la stessa
regola:

\begin{lstlisting}[language=Java]
// R5': isolamento SPI. Un'estensione tra moduli è una decisione esplicita:
// aggiungere una riga not(resideInAPackage(...)) con commento per ogni coppia sanzionata
// (ArchUnit 1.4 ha rimosso ignoreDependency dal fluent API di base).
@ArchTest
static final ArchRule does_not_depend_on_other_modules =
        noClasses()
                .that().resideInAPackage("..modules.offload..")
                .should().dependOnClassesThat(
                        resideInAPackage("..modules..")
                                .and(DescribedPredicate.not(resideInAPackage("..modules.offload.."))))
                .as("module must not depend on other modules (SPI isolation)");
\end{lstlisting}

Il divieto \`e sulle dipendenze \emph{laterali}. Dipendere dalle interfacce del
nucleo \`e il meccanismo previsto e resta libero: \`e cos\`i che i provider di
deployment implementano \code{ImageValidator} e che
\code{concurrency-control} scrive attraverso \code{ScalingMetricsSource}.

Il commento prescrive anche la procedura in caso di deroga: aggiungere una riga
esplicita con la propria giustificazione, cos\`i che ogni eccezione sia visibile
nel codice della regola. Al momento della stesura nessuna deroga \`e stata
aggiunta.

\subsection{Che cosa questo apparato compra}

Le sei regole insieme rendono \emph{eseguibile} l'architettura descritta nella
Parte II. Il \cref{cap:08-sistema-moduli} ha osservato che il compilatore \`e gi\`a
un esecutore parziale --- i moduli entrano come dipendenze
\code{runtimeOnly}, quindi il nucleo non pu\`o compilare contro di essi. Le
regole ArchUnit coprono ci\`o che il compilatore non vede: i cicli, la direzione
dei livelli, e la corrispondenza fra namespace e modulo fisico.

\section{Analisi statica}

\subsection{SonarQube}

L'analisi \`e su richiesta, non in integrazione continua:
\code{scripts/sonar.sh} avvia un container SonarQube effimero, analizza Java,
Python e Rust, e stampa il conteggio dei problemi aperti per linguaggio.

La rinuncia al \emph{quality gate} \`e argomentata esplicitamente:

\begin{quote}\itshape
Il server \`e effimero (nessun volume, nessuna storia), quindi SonarQube tratta
l'intera base di codice come ``nuovo codice'' a ogni esecuzione. Il gate di
default fallisce ogni volta che esiste un problema preesistente, il che
significa che un verdetto del gate sarebbe un FAIL costante e privo di
informazione.
\end{quote}

\`E la stessa logica che governa le soglie del gate di regressione
(\cref{cap:22-metodo-sperimentale}): un controllo che fallisce sempre viene
disattivato, e un controllo disattivato non controlla nulla. Meglio riportare un
numero che un verdetto inutile.

La conseguenza operativa \`e che la linea di base \`e zero problemi aperti, e
quindi \emph{qualunque} problema rilevato \`e nuovo. \`E un criterio pi\`u
severo di un gate sul nuovo codice, e sostenibile solo perch\'e la base di
codice \`e piccola.

\subsection{Analisi delle dipendenze}

Il progetto applica il plugin \code{dependency-analysis} a livello di settings,
che segnala dipendenze dichiarate e non usate, usate e non dichiarate, e
dipendenze che potrebbero cambiare configurazione.

Due osservazioni dall'esperienza registrata dal progetto: cinque dipendenze di
test sono state effettivamente rimosse a seguito dell'analisi, e gli starter di
Spring producono sistematicamente falsi positivi --- uno starter \`e per
costruzione una dipendenza che non si usa direttamente.

L'esclusione di \code{netty-codec-http3} descritta nel
\cref{cap:20-build-packaging} \`e il risultato pi\`u tangibile di questa linea di
lavoro: sei archivi rimossi, di cui cinque con codice nativo per piattaforma.

\subsection{Copertura}

JaCoCo su ogni sottoprogetto, con soglia minima all'85\%. I report XML
alimentano anche l'analisi SonarQube.

\section{Test di regressione strutturali}

La posizione del progetto sui test di prestazione \`e stata riportata nel
\cref{cap:21-osservabilita} e vale la pena richiamarla qui perch\'e appartiene a
questo capitolo: il repository contiene test \emph{strutturali} sul percorso
caldo anzich\'e microbenchmark assoluti. Asseriscono il riuso del replay, il
forward progress della coda sincrona dietro una testa bloccata, e l'equit\`a
asincrona fra funzioni calde e fredde.

\begin{quote}\itshape
Le temporizzazioni assolute sono sensibili all'ambiente, le garanzie di equit\`a
e riuso non lo sono.
\end{quote}

Ciascuno dei tre corrisponde a una propriet\`a discussa nella Parte III: il
replay al meccanismo di idempotenza (\cref{cap:06-nucleo}), il forward progress
alla scansione selettiva dello scheduler sincrono (\cref{cap:10-sync-queue}),
l'equit\`a al quanto di scheduling del modulo asincrono
(\cref{cap:09-async-queue}). Sono, in sostanza, i tre invarianti che i capitoli
precedenti hanno argomentato, resi eseguibili.

\section{Trappole note dell'impianto di test}

Il progetto ha accumulato alcune conoscenze operative che vale la pena
registrare, perch\'e ciascuna corrisponde a una propriet\`a reale del sistema.

\begin{description}[style=nextline, leftmargin=0pt]
  \item[\textbf{La coda sincrona respinge la prima richiesta.}] Senza storia di
  throughput la stima dell'attesa \`e infinita, e ogni richiesta viene respinta
  con \code{429} (\cref{cap:10-sync-queue}). I test di API devono quindi
  disattivare esplicitamente la coda sincrona. Non \`e un difetto: \`e il
  comportamento corretto in produzione che si manifesta in modo scomodo in
  laboratorio.

  \item[\textbf{Le metriche spariscono con la funzione.}] I meter per funzione
  sono rimossi alla cancellazione (\cref{cap:06-nucleo}), quindi una verifica
  che interroghi le metriche dopo aver cancellato la funzione non trova nulla.
  Gli scenari di validazione sono stati corretti mantenendo viva l'ultima
  funzione fino a dopo il controllo.

  \item[\textbf{Aggiungere un componente a un \code{record} rompe i test.}] Il
  costruttore canonico cambia firma e ogni istanziazione va aggiornata. \`E la
  ragione dei costruttori di compatibilit\`a descritti nel
  \cref{cap:05-contratti-dati}.

  \item[\textbf{Modificare i contratti condivisi richiede la suite intera.}] Una
  modifica a \code{platform/common} o a un SDK si propaga oltre il control plane;
  eseguire soltanto \code{:control-plane:test} non basta.
\end{description}

\section{Limiti dell'apparato}

\begin{enumerate}[leftmargin=*]
  \item \textbf{Lo spazio delle configurazioni modulari \`e coperto ai soli
  estremi} (\cref{cap:16-moduli-diagnostici}): \code{all} e \code{none}. Le
  configurazioni intermedie non hanno copertura dedicata.

  \item \textbf{L'analisi statica non \`e in integrazione continua.} \`E su
  richiesta, quindi la sua efficacia dipende dalla disciplina di chi sviluppa.

  \item \textbf{Gli scaffold generati non sono verificati end-to-end}
  (\cref{cap:19-cli-scaffolding}).
\end{enumerate}

I tre limiti hanno una struttura comune: riguardano tutti \emph{artefatti
generati o dichiarati} --- una configurazione di build, un file di
specifica, uno scaffold --- anzich\'e codice eseguito. \`E la categoria che
l'apparato di test copre peggio, ed \`e una lacuna riconoscibile e colmabile con
gli strumenti gi\`a presenti.
